\begin{corollary}[7.9.9]
\label{IV.2.s7.9.9-fr}
Sean $A$, $B$ dos anillos locales noetherianos y $\varphi:A\to B$ un
homomorfismo local que haga de $B$ un $A$-módulo plano. Supongamos que $A$ sea
íntegro y geométricamente unirrama, y que para toda $A$-álgebra finita
íntegra $A'$ se pueda resolver $\operatorname{Spec}(A')$ \textup{(lo que ocurrirá, por
ejemplo, si las fibras formales de $A$ son geométricamente regulares
y el cuerpo residual de $A$ tiene característica $0$, por los resultados de
Hironaka [35])}. Sea $k$ el cuerpo residual de $A$ y supongamos que
$\operatorname{Spec}(B\otimes_Ak)$ sea geométricamente íntegro en cada punto
\textup{(4.6.9)}. Entonces $B$ es íntegro.
\end{corollary}

\begin{proof}
Aplicaremos (7.9.8) tomando como $\boldsymbol{C}$ la categoría de
todos los preesquemas localmente noetherianos y como $\boldsymbol R$ la propiedad
de ser íntegro; los razonamientos de (7.5.3) y (7.9.7) muestran que las
condiciones $1^\circ$, $2^\circ$, $3^\circ$ y $4^\circ$ de (7.9.7) se
satisfacen en este caso. La hipótesis sobre $B\otimes_Ak$ y la proposición
(7.9.8) muestran, por tanto, (con las notaciones de (7.9.8)) que las fibras
$f^{-1}(y)$ son geométricamente íntegras en cada punto para $y\in Y$;
\emph{a fortiori}, los $f^{-1}(y)$ para $y\in Y$ son preesquemas reducidos,
y como $A$ es íntegro (y \emph{a fortiori} reducido), se ve ya que $B$
es \emph{reducido} (3.3.5). Queda por demostrar que $X=\operatorname{Spec}(B)$ es
\emph{irreducible}. Ahora bien, la demostración de (7.9.8) prueba que $X'$ es localmente
íntegro (por ser localmente noetheriano). Como el morfismo $g':X'\to X$ es
sobreyectivo, basta pues ver que $X'$ es irreducible o, puesto que $X'$
es localmente íntegro, que $X'$ es \emph{conexo}. Pero para ello basta
demostrar que la fibra $g'^{-1}(a)$ es \emph{conexa}. En efecto, establecido este punto,
$X'$ no puede ser suma de dos preesquemas abiertos no vacíos $X'_1$,
$X'_2$, pues se tendría entonces, por ejemplo,
$g'^{-1}(a)\cap X'_2=\varnothing$; pero como la restricción $X'_2\to X$ de $g'$
es propia (por ser $X'_2$ un subpreesquema cerrado de $X'$),
$g'(X'_2)$ sería una parte cerrada no vacía de $X$ que no contendría el punto
cerrado $a$, lo que es absurdo.

Ahora bien, se tiene
$g'^{-1}(a)=g^{-1}(b)\otimes_k \boldsymbol{k}(a)$; pero $g:Y'\to Y$ es propio y
birracional, por lo que, en la factorización de Stein
$Y'\xrightarrow{v}Y''\xrightarrow{u}Y$ del morfismo $g$ (\textbf{III}, 4.3.3),
el morfismo finito $u$ es también birracional y, por tanto, $Y''$ es íntegro.
Como se ha supuesto que $A$ es geométricamente unirrama, se deduce entonces de
(\textbf{III}, 4.3.4) que $g^{-1}(b)$ es \emph{geométricamente conexo}, con lo que
termina la demostración.
\end{proof}
